% ==============================================================================
% T0-Theorie / FFGFT: Dokument 185 (Kapitelinhalt)
% Titel: Die Zeit als Einbettungspreis der Mathematik:
%        Faktorisierung, Quantencomputer und die physikalische Grenze der RSA-Zerlegung
% Autor: Johann Pascher
% Datum: April 2026
% Bezüge:
%   Dok. 075: T0-Shor Algorithmus und Präzisionsbarriere
%   Dok. 076: Empirische Faktorisierungstests
%   Dok. 057: Relationales Zahlensystem, Primzahlen als irreduzible Verhältnisse
%   Dok. 183: Willow, T_rev, Rekursion und Zeit
% ==============================================================================

\section*{Abstract}

Faktorisierung ist in der reinen Mathematik eine zeitlose Identität: $N = p \cdot q$
existiert ohne Prozess, ohne Laufzeit, ohne Fehler. Sobald diese Identität jedoch
physikalisch realisiert wird — ob auf einem klassischen Prozessor, einem
Quantencomputer oder dem Vakuumfeld selbst — muss Information fließen. Fluss braucht
Zeit. Und Zeit, in der FFGFT definiert durch die fundamentale Umlaufzeit
$T_\text{rev} = \xipar \cdot \ell_P / c$, akkumuliert Phasenfehler.

Dieses Dokument zeigt: Nicht mangelnde Qubit-Zahl oder schlechte Hardware begrenzt
die RSA-Zerlegung durch Quantencomputer, sondern die intrinsische Zeit-Phasen-Kopplung
der physikalischen Raumzeit. Es gibt eine Skalierungsgrenze, jenseits derer selbst
der beste zukünftige Quantencomputer scheitert — geometrisch erzwungen, nicht
technisch umgehbar.

Gleichzeitig folgt aus derselben Geometrie ein positives Ergebnis: Fehlerkorrektur
kann massiv verbessert werden, wenn man berücksichtigt, dass das Vakuumfeld
bevorzugte Resonanzpunkte hat. Diese sind keine Artefakte — sie sind die stabilen
Torus-Moden, auf die das physikalische System natürlich zurückfällt.

\clearpage

\section{Das Paradoxon der Zeitlosigkeit}

\subsection{Zwei Welten: Mathematik und Physik}

In der reinen Mathematik ist die Faktorisierung $N = p \cdot q$ eine
\textbf{zeitlose Identität}. Es gibt keinen Prozess — nur eine statische Tatsache,
die gleichzeitig und immer wahr ist. Die Primzahl $p$ muss nicht gesucht werden;
sie ist bereits in $N$ enthalten, ohne jemals dort „hineingelegt" worden zu sein.

In der physikalischen Realität ist das anders. Ob auf Googles Willow-Prozessor,
einem klassischen Server oder im menschlichen Gehirn: Jeder Rechenvorgang erfordert,
dass Information von einem Ort zu einem anderen fließt. Fluss bedeutet Zeit.
Zeit bedeutet Phasendrift.

\begin{keypoint}[Die zentrale Unterscheidung]
\textbf{Mathematik:} $N = p \cdot q$ — zeitlos, fehlerfrei, ohne Prozess.

\textbf{Physik:} Um $p$ und $q$ aus $N$ zu \emph{extrahieren}, muss eine Messung
stattfinden — eine Periode $r$ muss in der Rückkopplung $f(x) = a^x \bmod N$
gefunden werden. Das erfordert Laufzeit $t = n \cdot T_\text{rev}$.
Laufzeit akkumuliert Phasenfehler. Ab einer kritischen Zahl $N$ überschreitet
der Fehler die Kohärenzschwelle, bevor die Resonanz gefunden ist.
\end{keypoint}

\subsection{Statisches Wissen vs. dynamischer Prozess}

Ein wichtiger Punkt: Bekannte Primzahlen müssen nicht erneut gesucht werden.
Sie existieren als bereits identifizierte, stabile Resonanzpunkte — eingetragen
in Tabellen, gespeichert in Datenbanken. Dieses \textbf{Wissen} ist statisch:
Es kostet keine Laufzeit, eine bekannte Primzahl abzurufen.

Aber die \emph{Verwendung} dieser bekannten Primzahlen zur Zerlegung einer neuen,
unbekannten Zahl $N$ ist kein statischer Abgleich — es ist ein \textbf{dynamischer
Prozess}. Das Einsetzen bekannter Primzahlen verkürzt zwar die Suche, aber die
Verknüpfung dieser Kandidaten mit $N$ erfordert eine physikalische Phasenkoinzidenz
zu prüfen, die Laufzeit benötigt.

\begin{technical}[Wissen vs. Prozess in FFGFT-Sprache]
In der Sprache der Torus-Resonanzen (Dok.~183):

\textbf{Bekannte Primzahl $p$:} Gespeicherte Windungszahl — ein topologisch
stabiler Knoten, der ohne weiteren Aufwand existiert.

\textbf{Faktorisierung von $N$:} Prüfung, ob die Windungszahl von $N$ als
Produkt zweier bekannter Windungszahlen $p \cdot q$ geschrieben werden kann.
Diese Prüfung erfordert, dass der Quantenprozessor die Phasenkoinzidenz zwischen
$p$, $q$ und $N$ physikalisch herstellt — in realer Zeit, mit akkumuliertem
Phasenfehler.
\end{technical}


\section{Faktorisierung als Resonanzsuche in der Zeit}

\subsection{Die Periodenfindung als physikalischer Prozess}

Shors Algorithmus (und der T0-Shor, Dok.~075) lösen Faktorisierung durch
\textbf{Periodenfindung}: Man sucht die Periode $r$ der Funktion
$f(x) = a^x \bmod N$. Diese Periode kodiert die Primfaktoren:
\begin{equation}
    a^r \equiv 1 \pmod{N}
    \;\Rightarrow\;
    a^r - 1 = (a^{r/2} - 1)(a^{r/2} + 1) \equiv 0 \pmod{N}
\end{equation}
Die gesuchte Resonanzfrequenz ist $\nu_r = 1/r$. Die benötigte Auflösung ist
(Dok.~075):
\begin{equation}
    \Delta f_\text{required} \approx \frac{f_0}{r^2}
\end{equation}
Für kryptographisch relevante Zahlen mit $r \approx N$:

\begin{table}[ht]
\centering
\begin{tabular}{|l|l|l|}
\hline
\textbf{Schlüssellänge} & \textbf{Benötigte Präzision $\epsilon$} & \textbf{Status} \\
\hline
27-Bit & $\epsilon \approx 10^{-16}$ & 64-Bit-Hardware ausreichend \\
\hline
56-Bit & $\epsilon \approx 10^{-34}$ & 128-Bit-Hardware ausreichend \\
\hline
RSA-1024 & $\epsilon \approx 10^{-617}$ & physikalisch unmöglich \\
\hline
RSA-2048 & $\epsilon \approx 10^{-1234}$ & physikalisch unmöglich \\
\hline
\end{tabular}
\caption{Präzisionsbarriere für Faktorisierung (nach Dok.~075)}
\end{table}

\subsection{Warum Laufzeit-Akkumulation unausweichlich ist}

In FFGFT ist Zeit keine externe Uhr, sondern die akkumulierte Laufzeit des Feldes
durch seine eigene Geometrie — gemessen in Vielfachen von $T_\text{rev}$ (Dok.~183):
\begin{equation}
    T_\text{rev} = \frac{\xipar \cdot \ell_P}{c} \approx 7{,}2 \times 10^{-48}\,\text{s}
\end{equation}
Um die Periode $r$ zu finden, muss das System $r$ Umläufe vollziehen — also eine
Gesamtlaufzeit von $t = r \cdot T_\text{rev}$. Bei jedem Umlauf akkumuliert ein
Phasenfehler $\varepsilon$. Nach $r$ Zyklen beträgt der Gesamtfehler $r \cdot \varepsilon$.

Die Kohärenz bricht zusammen, sobald der Phasenfehler $\pi/2$ überschreitet:
\begin{equation}
    r_\text{krit} \approx \frac{\pi}{2\varepsilon}
\end{equation}
Für RSA-Schlüssel ist $r \approx N$ gigantisch. Die benötigte Präzision pro Schritt
skaliert als $\varepsilon \lesssim \pi/(2N)$ — und das erfordert für RSA-1024
die physikalisch unmögliche Präzision $\epsilon \approx 10^{-617}$.

\begin{foundation}[Die Laufzeit-Akkumulation als Naturgesetz]
Das ist keine technische Einschränkung, die durch bessere Hardware überwunden
werden kann. Es ist eine Konsequenz der Struktur der physikalischen Realität:

\begin{equation}
    \Psi(t) = G\!\left(\Psi(t - T_\text{rev}),\; \xipar\right)
\end{equation}

Jeder Rechenschritt ist ein Umlauf auf dem Torus. Jeder Umlauf akkumuliert
$\Delta\phi = 2\pi \cdot \xipar$. Für $r \approx N$ Umläufe bei RSA-Schlüsseln
ist der akkumulierte Fehler $r \cdot \Delta\phi \approx N \cdot 2\pi \cdot \xipar$
— weit jenseits der Kohärenzschwelle.
\end{foundation}


\section{Warum der Quantencomputer dennoch scheitert}

\subsection{Das Threshold-Theorem und seine Grenzen}

Die Standard-Quanteninformatik verweist auf das \textbf{Threshold-Theorem}:
Wenn die Fehlerrate pro Gate unter einem kritischen Schwellenwert liegt, kann
Quantum Error Correction (QEC) Fehler schneller korrigieren als sie entstehen.
Der logische Zustand lebt dann theoretisch beliebig lang.

Die FFGFT zeigt eine fundamentale Grenze dieses Arguments:

\begin{important}[QEC ist selbst ein rekursiver Prozess]
QEC benötigt Laufzeit — jede Fehlerkorrektur ist selbst eine Sequenz von
Gate-Operationen, die $T_\text{rev}$ kostet. Die Korrektur fügt dem System
neue Phasenfehler hinzu, die ihrerseits korrigiert werden müssen.

Für kleine $N$ (kurze Perioden, wenige Zyklen) überwiegt die Korrektur den
Fehler — das Threshold-Theorem gilt. Für große $N$ (RSA-Schlüssel) wird die
Korrektur selbst Teil des Problems: Die Rekursionstiefe der Fehlerkorrektur
ist so groß, dass der Phasenfehler der Korrektur denselben Charakter annimmt
wie das ursprüngliche Rauschen.
\end{important}

Die Argumentation der Standardliteratur — „Fehlerkorrektur besiegt die Zeit" —
gilt im Regime kurzer Perioden. Im RSA-Regime mit $r \approx 10^{308}$ (RSA-1024)
versagt sie prinzipiell.

\subsection{RSA als geometrischer Schutzschild}

In der FFGFT-Perspektive ist die Schwierigkeit der RSA-Faktorisierung kein
Mangel an Rechenleistung, sondern ein \textbf{Schutzschild der Raumzeit-Geometrie}:

Große Primzahlen erzeugen im Rückkopplungssystem $a^x \bmod N$ so
langperiodische Resonanzen, dass kein physikalisches System — das selbst Teil
dieser Raumzeit ist und selbst durch $T_\text{rev}$ und $\xipar$ beschrieben wird
— die nötige Phasenstabilität über die Laufzeit $r \cdot T_\text{rev}$ aufrecht
erhalten kann.

\begin{conclusionBox}[Die Mauer zwischen Mathematik und Physik]
Die Mathematik sagt: $N = p \cdot q$ ist trivial wahr.

Die Physik sagt: Um diese Wahrheit zu \emph{entdecken}, muss ein System eine
Periode $r \approx N$ Zyklen kohärent durchhalten. Für RSA-1024 bedeutet das
$r \approx 10^{308}$ Zyklen mit Präzision $\epsilon \approx 10^{-617}$ pro Zyklus.

Das ist die Mauer. Nicht aus Metall — aus Zeit und Phase.
\end{conclusionBox}


\section{Das positive Ergebnis: Resonanzpunkte verbessern die Fehlerkorrektur}

\subsection{Bevorzugte Resonanzpunkte im Torus-Spektrum}

Aus derselben Geometrie, die die RSA-Zerlegung begrenzt, folgt ein positives
Ergebnis für die Quantenfehlerkorrektur. Das Torus-Spektrum ist nicht gleichmäßig:
Die Moden tragen mit $1/n^2$ bei, und Primzahlmoden sind die irreduziblen
Resonanzknoten. Das bedeutet: Es gibt \textbf{bevorzugte, geometrisch stabile
Positionen} im Zustandsraum des Quantencomputers.

Die Standardkorrektur (Threshold-Theorem) behandelt alle Fehler statistisch
gleichwertig und projiziert auf einen beliebigen Zielzustand. Die FFGFT-informierte
Fehlerkorrektur erkennt, dass bestimmte Zustände \textbf{geometrisch bevorzugt}
sind — sie entsprechen stabilen Torus-Moden, auf die das System natürlich
zurückfällt.

\subsection{Resonanzkorrektur statt statistischer Korrektur}

Statt jeden Fehler unabhängig zu korrigieren, kann man Fehler als
\textbf{Abdriften von stabilen Resonanzpunkten} interpretieren. Die Korrektur
besteht dann darin, das System aktiv auf den nächsten stabilen Resonanzknoten
zurückzuführen — nicht auf einen beliebigen Codewort-Zustand.

\begin{keyresult}[Verbesserungspotenzial durch Resonanzkorrektur]
Wenn die Fehlerkorrektur geometrisch informiert ist — also die bevorzugten
Torus-Moden als Ankerpunkte nutzt — kann sie Fehler effizienter korrigieren als
die statistische Methode. Die Rate, mit der Fehler akkumulieren, bleibt dieselbe;
aber die Rate, mit der sie korrigiert werden, steigt, weil die Korrekturbewegungen
auf natürliche Attraktoren des Systems zeigen.

Dies ist kein Widerspruch zur fundamentalen Laufzeit-Barriere: Für kleine $N$
hilft Resonanzkorrektur erheblich. Für RSA-Schlüssellängen übersteigt die
Periode $r$ die Reichweite jeder Korrektur.
\end{keyresult}

\subsection{Die praktische Konsequenz für Willow und Nachfolger}

Googles Willow demonstriert exponentielle Fehlerunterdrückung mit wachsender
Code-Distanz $d$ (Dok.~183) — dies ist der experimentelle Beleg, dass topologisch
stabile Moden zugänglich sind. Die FFGFT-Perspektive sagt:

\begin{itemize}
    \item \textbf{Was Willow bereits tut:} Projektion auf topologisch stabile
    Torus-Moden — implizit, ohne die geometrische Struktur zu kennen.
    \item \textbf{Was durch Resonanzkorrektur möglich wäre:} Explizite Nutzung
    der $1/n^2$-Struktur des Torus-Spektrums, um Korrekturen gezielt auf die
    stabilsten Moden zu richten.
    \item \textbf{Was prinzipiell unmöglich bleibt:} RSA-Schlüssel im kryptographisch
    relevanten Bereich zu brechen — weil die Periode $r \approx N$ die Kohärenzzeit
    jeder physikalischen Hardware übersteigt.
\end{itemize}


\section{Gegenargumente und FFGFT-Antworten}

\begin{table}[ht]
\centering
\begin{tabular}{|p{5cm}|p{7.5cm}|}
\hline
\textbf{Standard-Argument} & \textbf{FFGFT-Antwort} \\
\hline
Threshold-Theorem: QEC besiegt die Zeit & QEC ist selbst ein rekursiver Prozess
mit Laufzeit. Für $r \approx 10^{308}$ wird die Korrektur Teil des Rauschens. \\
\hline
Shor ist polynomial: $O((\log N)^3)$ Schritte & Die Polynomialität gilt für
die logische Komplexität. Die Phasenpräzision $\epsilon \sim 1/N^2$ wächst
exponentiell — physikalisch, nicht logisch. \\
\hline
Bessere Hardware $\Rightarrow$ kleineres $\epsilon$ & Es gibt kein $\epsilon = 0$.
Jede Hardware ist in die Raumzeit eingebettet, also durch $T_\text{rev}$ und
$\xipar$ beschrieben. Die geometrische Phasendrift ist intrinsisch. \\
\hline
Topologische Qubits lösen das Problem & Topologische Qubits verbessern
die Fehlerrate — sie beseitigen die Laufzeit-Akkumulation nicht. Für RSA-1024
reicht auch $\epsilon \to 10^{-16}$ nicht aus: man bräuchte $10^{-617}$. \\
\hline
\end{tabular}
\caption{Standard-Argumente der Quanteninformatik und FFGFT-Antworten}
\end{table}


\section{Zusammenfassung}

Die Zeit ist der Zoll, den man zahlt, wenn man Mathematik in die physikalische
Realität bringt. Für kleine Zahlen ist dieser Zoll erschwinglich — Quantencomputer
können und werden die Kryptographie kurzer Schlüssel brechen. Für RSA-Schlüssel
der gegenwärtigen und zukünftigen Kryptographie ist der Zoll astronomisch.

\begin{foundation}[Drei Aussagen dieses Dokuments]
\begin{enumerate}
    \item \textbf{Physikalische Grenze:} RSA-1024 und größer bleibt durch
    die Laufzeit-Akkumulation der Raumzeit-Geometrie ($T_\text{rev}$, $\xipar$)
    vor Quantencomputern geschützt — nicht durch technische Mängel, sondern
    durch Naturgesetze.

    \item \textbf{Statisches Wissen vs. dynamischer Prozess:} Bekannte Primzahlen
    kosten keine Laufzeit. Ihre Anwendung auf eine neue Zahl $N$ erfordert
    eine physikalische Phasenkoinzidenz — und die kostet $r \cdot T_\text{rev}$.

    \item \textbf{Positives Ergebnis:} Fehlerkorrektur kann durch geometrisch
    informierte Resonanzkorrektur (Nutzung bevorzugter Torus-Moden) erheblich
    verbessert werden — für kleine $N$, für kurze Perioden, für praktische
    Quantenalgorithmen im erreichbaren Bereich.
\end{enumerate}
\end{foundation}

\vspace{0.8cm}
\noindent
\textit{Entwurf. Die Aussage zur physikalischen Unmöglichkeit der RSA-Zerlegung
bei großen Schlüsseln ist qualitativ aus der FFGFT-Dynamik motiviert. Eine
quantitative Herleitung der genauen Skalierungsgrenze erfordert weitere Formalisierung
der Verbindung zwischen $T_\text{rev}$, $\xipar$ und der Threshold-Bedingung des
Shor-Algorithmus.}

\begin{thebibliography}{9}

\bibitem{Dok075}
J.~Pascher, \textit{T0-Shor Algorithmus: Energiefeld-Formulierung und
Präzisionsbarriere}, T0-Dokument~075 (2025).

\bibitem{Dok076}
J.~Pascher, \textit{Empirische Faktorisierungstests}, T0-Dokument~076 (2025).

\bibitem{Dok057}
J.~Pascher, \textit{Relationales Zahlensystem: Primzahlen als irreduzible
Verhältnisse}, T0-Dokument~057 (2025).

\bibitem{Dok183}
J.~Pascher, \textit{Google Willow und die Physik der Selbstmessung: $T_\text{rev}$
als fundamentale Umlaufzeit}, T0-Dokument~183 (2026).

\bibitem{Shor1997}
P.~W.~Shor, \textit{Polynomial-Time Algorithms for Prime Factorization and
Discrete Logarithms on a Quantum Computer}, SIAM J.\ Comput.\ \textbf{26}
(1997), 1484--1509.

\bibitem{GoogleWillow2024}
Google Quantum AI, \textit{Quantum error correction below the surface code
threshold}, Nature (2024), DOI: 10.1038/s41586-024-08449-y.

\end{thebibliography}
